\begin{activity}{Solving the Heat Equation as a Human Computer}

\section*{Purpose}
To develop intuition for finite difference methods by acting as the nodes in a
numerical heat flow model.

In a numerical model, each point (node) updates its temperature using only
information from its neighbours.

\emph{In this activity, you will act as the nodes in a finite difference model.}

{\color{red}
\section*{Learning Outcomes}

By the end of this activity, you should be able to:
\begin{itemize}
  \item Apply the finite difference update rule by computing temperature from neighbouring values.
  \item Predict qualitatively how an initial temperature distribution evolves under the diffusion update rule.
  \item Explain why temperature peaks smooth out and why the system approaches a uniform (or linear) steady state.
  \item Describe why heat transfer is a local process and how information about temperature propagates through a grid one node at a time.
\end{itemize}
}

\ntextbox{Human Computer}{
    Before electronic computers existed, a \emph{computer} was literally a person employed to carry out calculations by hand or with mechanical calculators. Large groups of human computers were used in astronomy, navigation, ballistics, engineering, and later for the numerical solution of differential equations. By the late 19th and early 20th centuries, many of these computing groups were composed largely of women. Examples include the astronomical computers working for the \emph{Nautical Almanac} and later teams at observatories, government laboratories, and research institutions.

    During World War II, the scale of numerical work \emph{exploded} (pun intended). Hundreds of human computers computed ballistic firing tables and solved systems of finite-difference equations for artillery trajectories, fluid flow, and other military problems. Post-WWII, human computers were also instrumental to the US space program. If you've seen the movie \textit{Hidden Figures}, you learned about some of these remarkable women who helped make going to the Moon possible by producing many of the calculations needed to determine flight plans, rocket burn parameters, launch windows, and safe re-entry trajectories.

    Finite-difference methods were especially important because many physical problems could not be solved analytically. Human computers discretized differential equations onto grids and repeatedly computed numerical updates row by row—an extremely laborious process that helped motivate the development of electronic computers. The first six programmers of the world's first general-purpose electronic computer, ENIAC, were women who had previously worked as human computers.

    David Alan Grier, \textbf{2005}, \textit{When Computers Were Human}, Princeton University Press. \\
    Kathy Kleiman, \textbf{2022}, \textit{Proving Ground: The Untold Story of the Six Women Who Programmed the World's First Modern Computer}. \\
    David Alan Grier, \textbf{2001}, \textit{Human computers: the first pioneers of the information age}, \textbf{Endeavour} 25(1).
}{aside}

\section*{Setup}

\begin{itemize}
    \item Each student is assigned a position in a line.
    \item Each student is given an initial temperature value $T_i$.
    \item The two students at the ends of the line represent \textbf{fixed
    boundary conditions}.
\end{itemize}

\section*{Update rule}
At each step, interior nodes update their temperature according to:
\begin{equation}
    T_i^{\text{new}} = \frac{T_{i-1} + T_{i+1}}{2}
\end{equation}

This means:
\begin{itemize}
    \item you look at your two neighbours,
    \item take the average,
    \item and update your value.
\end{itemize}

\begin{questions}
    \question{Prediction}
    Before starting, discuss:
    \begin{itemize}
        \item What happens to a sharp temperature spike over time?
        \item Will the system become smoother or less smooth?
    \end{itemize}
    \answerbox[3cm]

    \question{Iteration}
    Repeat the following for several steps:
    \begin{enumerate}
        \item Show your value to your neighbours.
        \item Compute your new value.
        \item Update simultaneously.
    \end{enumerate}

    \textbf{Record the evolution of the system.} (Use the table below to track temperature values at each node across three time steps.)

    \begin{center}
        \begin{tblr}{
            colspec = {c*{9}{X[c]}},
            row{2-Z} = {ht=2.5em},
            hline{1,2,Z} = {solid},
            hline{3-{Z-1}} = {dashed},
            vline{1,2,11} = {solid},
            vline{3-10} = {dashed}
        }
            \diagbox{\textbf{Node}}{\textbf{Step}} & \textbf{1} & \textbf{2} & \textbf{3} & \textbf{4} & \textbf{5} & \textbf{6} & \textbf{7} & \textbf{8} & \textbf{9} & \textbf{10} \\
            0 & & & & & & & & & \\
            1 & & & & & & & & & \\
            2 & & & & & & & & & \\
            3 & & & & & & & & & \\
            4 & & & & & & & & & \\
            5 & & & & & & & & & \\
            6 & & & & & & & & & \\
            7 & & & & & & & & & \\
            8 & & & & & & & & & \\
        \end{tblr}
    \end{center}

    \question{Observations}
    Discuss:
    \begin{itemize}
        \item Where do the largest changes occur?
        \item What happens to temperature peaks?
        \item How does heat move through the system?
    \end{itemize}
    \answerbox[3cm]

    \question{Interpretation}

    \textbf{Explain in words:}
    \begin{itemize}
        \item Why does the system become smoother over time?
        \item What role do neighbouring nodes play?
        \item How does this relate to diffusion?
    \end{itemize}
    \answerbox[4cm]

\end{questions}

\bigskip
\hrule
\bigskip

\section*{Optional Extension 1: Heat Production}

Introduce a constant heat source at one node.

\textbf{New update rule:}
\begin{equation}
T_i^{\text{new}} = \frac{T_{i-1} + T_{i+1}}{2} + S
\end{equation}
where $S$ is a constant heat input.

\textbf{Tasks:}
\begin{itemize}
    \item What happens to the system over time?
    \item Does it still become uniform?
    \item What steady-state behaviour emerges?
\end{itemize}

\bigskip
\hrule
\bigskip

\section*{Optional Extension 2: Two-Dimensional Grid}

Arrange students in a grid.

\textbf{Update rule:}
\begin{equation}
    T_{i,j}^{\text{new}} =
    \frac{
    T_{i+1,j} + T_{i-1,j} + T_{i,j+1} + T_{i,j-1}
    }{4}
\end{equation}


\textbf{Tasks:}
\begin{itemize}
    \item Observe how a hot region spreads in two dimensions.
    \item Compare to the one-dimensional case.
    \item Describe the shape of the temperature field over time.
\end{itemize}

\section*{Final Reflection}

\begin{itemize}
    \item In what sense are you “solving the heat equation”?
    \item Why does this method require many iterations?
\end{itemize}


\section*{Demo}

There is a simple python GUI, \verb+finitediff_gui.py+, that you can use to explore calculations similar to the ones we just did in this exercise.  The implementation solves the Poisson equation:
\[
    \laplacian \T = - \frac{A}{k}
\]
using the discrete equation
\[
    T_i^{n+1} = 
    \frac{k_{i-1}T_{i-1}^n+k_{i+1}T_{i+1}^n} {k_{i-1}+k_{i+1}} 
    + \frac{A_i}{k_i}
\]


\section*{Key Ideas}

Heat flow is a local process: each point evolves based on its neighbours.

\textbf{Numerical models evolve through local interactions.} Each node only communicates with its immediate neighbours. Global behaviour (e.g., a uniform steady state) emerges from many local interactions over many time steps. There is no global solution found in one step.

\textbf{Information propagates at a finite rate.} Because each update depends only on nearest neighbours, a temperature change at one end of the chain cannot instantaneously affect the other end. It must propagate node by node, step by step. This is the discrete analogue of finite diffusion speed in the continuum.

\textbf{Steady state is reached when updates become zero.} The system reaches a steady state when every node's temperature equals the average of its neighbours: $T_i = (T_{i-1} + T_{i+1})/2$. This is the discrete form of $\nabla^2 T = 0$ --- the Laplace equation for heat without sources.

\end{activity}